% Generated from audited fresh results. Do not edit numeric content.
\begin{table*}[t]
\centering
\setlength{\tabcolsep}{4pt}
\begin{tabular}{lrrrrr}
\toprule
Policy & e-SNLI-VE & SugarCrepe & SC++ both & COCO I$\to$T & COCO T$\to$I\\
\midrule
Frozen & 79.47 & 81.21 & 68.95 & 56.32 & 39.35\\
\midrule
Source positives & $81.18_{\pm 0.14}$ & $83.26_{\pm 0.30}$ & 67.40 & 55.85 & 39.38\\
Positive expansion & $87.00_{\pm 0.40}$ & $83.39_{\pm 0.15}$ & 67.23 & 47.29 & 34.84\\
+ contradiction weighting & $87.28_{\pm 0.33}$ & $83.33_{\pm 0.16}$ & 67.30 & 47.25 & 34.81\\
+ neutral exclusion & $86.66_{\pm 0.44}$ & $83.39_{\pm 0.27}$ & 67.66 & 48.01 & 35.39\\
+ both & $86.74_{\pm 0.17}$ & $83.24_{\pm 0.31}$ & 67.61 & 47.26 & 35.27\\
Constant weights & $82.15_{\pm 0.24}$ & $83.67_{\pm 0.20}$ & 68.37 & 55.21 & 39.39\\
\bottomrule
\end{tabular}
\caption{Relation interventions with the same candidate pool and tuning budget. Scores are percentages; subscripts show sample standard deviation across three selected seeds for the primary endpoints. COCO columns report full-pool R@1. The three rows marked + add rules to positive expansion. Constant weights use source positives.}
\label{tab:main}
\end{table*}

\paragraph{The largest benefit is on the annotated distinction.}
Positive expansion improves e-SNLI-VE by \PositiveRelationGain{} percentage points over source positives. Its SugarCrepe gain is \PositiveSugarGain{} points. Table~\ref{tab:main} separates the two negative-pair interventions on top of that expanded set. The combined policy differs from positive expansion by \CombinedRelationDelta{} points on e-SNLI-VE and \CombinedSugarDelta{} on SugarCrepe. Their adjusted 99.17\% image-bootstrap intervals are $\CombinedRelationInterval$ and $\CombinedSugarInterval$.

\paragraph{The edited-caption gain targets added content.}
SugarCrepe mixes edits that add objects or attributes with edits that replace content or swap its assignment. The aggregate score hides opposite effects. Positive expansion improves the added-content group and reduces accuracy on the combined replacement/swap group (Table~\ref{tab:edits}). Both directions hold in every seed. Appendix~\ref{app:results} reports the seven original categories.
% Generated from audited fresh results. Do not edit numeric content.
\begin{table}[t]
\centering
\setlength{\tabcolsep}{3pt}
\begin{tabular}{lrrrr}
\toprule
Edit group & $n$ & Source & Expand. & $\Delta$\\
\midrule
Added content & 2,754 & 88.09 & 90.01 & +1.92\\
Replace/swap & 4,757 & 80.46 & 79.56 & -0.90\\
All edits & 7,511 & 83.26 & 83.39 & +0.14\\
\bottomrule
\end{tabular}
\caption{SugarCrepe accuracy by edit group. Differences are positive expansion minus source positives, in percentage points. Groups partition the official pairs.}
\label{tab:edits}
\end{table}


\paragraph{Full-pool retrieval exposes the practical cost.}
Positive expansion loses \PositiveRetrievalLoss{} points of image-to-text R@1 relative to source positives. On the fixed pool, that is \ExtraFailedImageQueries{} fewer image queries returning an annotated correct caption first per 5,000 queries, averaged across seeds. Every seed loses retrieval accuracy. Text-to-image R@1 falls by \PositiveTextRetrievalLoss{} points as well. These tests keep all candidate captions and images in the pool, so the loss concerns the rank of a usable match rather than a binary choice against one designed foil.
% Generated from audited fresh results. Do not edit numeric content.
\begin{figure}[t]
\centering\includegraphics[width=0.78\columnwidth]{retrieval_query.jpg}
\par\raggedright\textbf{Source positives:} A young girl is preparing to blow out her candle.

\textbf{Positive expansion:} a woman is making some food in a bowl
\caption{Top retrieved captions for the same image query (seed 17). The best annotated caption moves from rank 1 to 6. Both models search the full caption pool.}\label{fig:query}
\end{figure}


Figure~\ref{fig:query} shows a changed top result for the same image and caption pool. A lower score for a deliberately edited foil would not detect this change in the match returned to the user.

\paragraph{A simple control preserves a different set of capabilities.}
Constant negative weighting retains \ConstantExpandedRetrievalGain{} points more image-to-text R@1 than positive expansion. It also scores higher on edited captions, while its annotated statement accuracy is lower (Table~\ref{tab:main}). This control weakens all negative competition without deciding which pairs are contradicted or neutral. The selected means expose the application choice: improving the annotation distinction is different from preserving edited-caption and retrieval performance. Table~\ref{tab:weights} later compares the similarity-based weighting rules.

\paragraph{The use of neutral labels has a conditional benefit.}
Neutral exclusion beats count-matched random removal by \NeutralRandomRelationGain{} points on e-SNLI-VE and \NeutralRandomSugarGain{} on SugarCrepe. However, when contradiction weighting is retained, adding neutral exclusion changes e-SNLI-VE by \NeutralUnderHardeningDelta{} points, with adjusted interval $\NeutralUnderHardeningInterval$. Selecting neutral pairs is therefore useful relative to random removal, but it does not improve the combined selected procedure over contradiction weighting alone.

\paragraph{Alternative wordings test a separate requirement.}
All \NonzeroAdapterCount{} nonzero-epoch adapters reduce accuracy on the alternative valid SugarCrepe++ caption relative to the frozen model. This occurs across policies, including source positives. Table~\ref{tab:main} reports the stricter score that requires both valid captions to beat the foil. Appendix~\ref{app:retention} shows the paired first-caption and alternative-caption outcomes for every selected adapter.
